\documentclass[conference]{IEEEtran}
\usepackage{cite}
\usepackage{amsmath,amssymb,amsfonts}
\usepackage{algorithmic}
\usepackage{graphicx}
\usepackage{textcomp}
\usepackage{xcolor}
\usepackage{url}

\def\BibTeX{{\rm B\kern-.05em{\sc i\kern-.025em b}\kern-.08em
    T\kern-.1667em\lower.7ex\hbox{E}\kern-.125emX}}

% Unresolved placeholders are shown in red. Search for \todo to find them all.
\newcommand{\todo}[1]{\textcolor{red}{[#1]}}

\begin{document}

\title{The Impact of Generative AI Coding Assistants on Developer Code Comprehension in Software Development and Maintenance: A Systematic Literature Review}
% Note: Sub-titles are not captured in Xplore and should not be used.
% CRITICAL: Do not use symbols, special characters, footnotes, or math in
% the paper title or abstract.

\author{
\IEEEauthorblockN{1\textsuperscript{st} Athalla Barka Fadhil}
\IEEEauthorblockA{\textit{Department of Information Technology} \\
\textit{Institut Teknologi Sepuluh Nopember}\\
Surabaya, Indonesia \\
5027231018@student.its.ac.id}
\and
\IEEEauthorblockN{2\textsuperscript{nd} Naufal Syafi' Hakim}
\IEEEauthorblockA{\textit{Department of Information Technology} \\
\textit{Institut Teknologi Sepuluh Nopember}\\
Surabaya, Indonesia \\
5027231022@student.its.ac.id}
\and
\IEEEauthorblockN{3\textsuperscript{rd} Muhammad Dzaky Ahnaf}
\IEEEauthorblockA{\textit{Department of Information Technology} \\
\textit{Institut Teknologi Sepuluh Nopember}\\
Surabaya, Indonesia \\
5027231039@student.its.ac.id}
\and
\IEEEauthorblockN{4\textsuperscript{th} Dimas Andhika Diputra}
\IEEEauthorblockA{\textit{Department of Information Technology} \\
\textit{Institut Teknologi Sepuluh Nopember}\\
Surabaya, Indonesia \\
5027231074@student.its.ac.id}
}

\maketitle

\begin{abstract}
\todo{Write the abstract last, after the final study selection, synthesis, discussion, and conclusions are verified.}
\end{abstract}

\begin{IEEEkeywords}
generative AI, coding assistants, large language models, code comprehension, software maintenance, systematic literature review
\end{IEEEkeywords}

%% ====================================================================
%% I. INTRODUCTION
%% ====================================================================
\section{Introduction}\label{sec:intro}

Understanding existing source code is one of the most frequent activities in software engineering. In a field study of 78 professional developers working on seven real projects, Xia et al.\ found that developers spent roughly 58\% of their working time on program comprehension activities \cite{xia2018}. Software maintenance, in which engineers must modify, debug, and review code they did not write, depends heavily on this activity.

Generative AI (GenAI) coding assistants built on large language models (LLMs) have rapidly entered this workflow. The Codex model introduced in 2021, whose production version powers GitHub Copilot \cite{chen2021}, marked the arrival of developer-facing assistants, and adoption has since become mainstream: in the 2025 Stack Overflow Developer Survey, 84\% of respondents reported using or planning to use AI tools in their development process, up from 76\% a year earlier \cite{so2025}. These assistants change what comprehension demands of a developer. Developers can now ask an assistant to explain unfamiliar code, but they must also understand, validate, and maintain code that an assistant produced and that no human colleague wrote.

Recent empirical work reflects both sides of this change. On the supportive side, researchers have designed and evaluated assistants intended to help developers understand code, including LLM-based explanations inside the IDE \cite{nam2024}, anchored explanations of just-generated code \cite{yan2024}, LLM-generated context that enriches code comprehension \cite{grandel2026}, and interactive visual aids combined with LLMs for navigating unfamiliar codebases \cite{kim2026}. On the cautionary side, studies of how developers validate and repair LLM-generated code examine their behaviour using IDE traces, eye tracking, and cognitive-workload measures \cite{tang2024}, or test whether live programming helps them validate generated programs \cite{ferdowsi2024}. A study of proactive AI programming support reports trade-offs between assistance and developers' understanding, including comprehension loss and cognitive burden \cite{pu2025}. Practitioners voice a similar concern: the most common frustration in the same survey, cited by 66\% of respondents, is AI solutions that are ``almost right, but not quite'' \cite{so2025}, which are outputs a developer has to read and check before trusting.

Existing secondary studies address adjacent questions. Mohamed et al.\ reviewed 39 studies on how LLM assistants affect developer productivity, and note risks such as cognitive offloading \cite{mohamed2026}. Qiao et al.\ reviewed 31 studies published between 2022 and 2024 on GenAI assistants for code comprehension, but frame their implications for computing education \cite{qiao2026}. To the best of our knowledge, no systematic review has yet synthesized the evidence on how GenAI coding assistants affect the code comprehension of developers performing realistic software development and maintenance tasks. This evidence is also hard to consolidate by reading individual studies. The primary studies cited above appear in software engineering, human--computer interaction, and program comprehension venues, and they operationalize comprehension in different ways, for example through comprehension quizzes and cognitive-load instruments \cite{kim2026} or through eye tracking \cite{tang2024}. Productivity, task completion time, or code quality alone do not show whether a developer understands the code, so studies that report only those outcomes cannot answer the question either.

To address this gap, this paper presents a systematic literature review (SLR) conducted according to the guidelines of Kitchenham and Charters \cite{kitchenham2007}, with study flow reported in the style of PRISMA 2020 \cite{page2021}. The review covers peer-reviewed primary empirical studies published between January 2021 and December 2026 in which human developers use interactive GenAI or LLM-based coding assistants for realistic software development or maintenance work, and in which code comprehension is measured directly or through a close proxy such as cognitive load, knowledge retention, or the ability to correctly modify or debug code. Model-only benchmark studies and studies confined to introductory programming education are out of scope. The review is guided by four research questions:

\begin{itemize}
\item \textbf{RQ1:} How does the use of GenAI coding assistants affect developers' code comprehension during software development and maintenance tasks?
\item \textbf{RQ2:} What GenAI coding assistants and forms of comprehension support have been used by developers in these studies?
\item \textbf{RQ3:} How has developer code comprehension been measured, and in what study designs, tasks, and codebase contexts?
\item \textbf{RQ4:} What challenges limit developers' code comprehension when working with GenAI assistants?
\end{itemize}

This paper makes the following contributions:
\begin{enumerate}
\item A documented, protocol-driven evidence base of primary empirical studies, with eligibility criteria that separate evidence of developer comprehension from evidence of productivity, code quality, or model accuracy alone.
\item A catalogue of the GenAI coding assistants, interaction modes, and forms of comprehension support used in these studies (RQ2).
\item An overview of how comprehension has been measured and in which study designs, tasks, and codebase contexts, distinguishing direct measures from close proxies (RQ3).
\item A synthesis of the reported effects on comprehension and of the challenges that limit it (RQ1 and RQ4), with implications for tool designers, practitioners, and future studies.
\end{enumerate}

The remainder of this paper is organized as follows. Section~\ref{sec:method} describes the review methodology, including the search strategy, selection criteria, quality assessment, and data extraction. Section~\ref{sec:results} presents the results for each research question. Section~\ref{sec:discussion} discusses the findings and their implications, Section~\ref{sec:threats} discusses threats to validity, and Section~\ref{sec:conclusion} concludes the paper.

%% ====================================================================
%% PLACEHOLDER SECTIONS (replace with the next chapters)
%% ====================================================================
\section{Methodology}\label{sec:method}

This systematic literature review (SLR) was designed using the evidence-based software engineering approach of Kitchenham and Charters \cite{kitchenham2007}, with reporting informed by SEGRESS \cite{segress2023}, PRISMA 2020 \cite{page2021}, and PRISMA-S \cite{prismaS2021}. A predefined protocol specified the research questions in Section~\ref{sec:intro}, eligibility criteria, search sources, quality appraisal, and planned synthesis. The PICOC framework delimited human developers (including advanced students performing realistic engineering tasks), interactive GenAI coding assistance, optional comparator conditions, developer code comprehension as the outcome, and software development or maintenance as the context. Departures from the planned search were recorded in a protocol-change log.

\subsection{Search Strategy}
Scopus, IEEE Xplore, and ScienceDirect were searched on 17 September 2026 for peer-reviewed journal and conference studies published from 2021 onward, with 2026 as the upper publication-year limit and English as the language limit. Search terms combined three concepts with Boolean AND: GenAI or LLM-based coding assistance; code comprehension and related outcomes; and programming or software. Synonyms and selected assistant names within each concept were joined by OR. Scopus used \texttt{TITLE-ABS-KEY}; IEEE Xplore used \texttt{All Metadata}; and ScienceDirect used the title, abstract, or author-specified keywords field. ScienceDirect required five narrower query variants. One complete executed ScienceDirect query (SD-S1) was:

\begin{quote}
\footnotesize\raggedright
\texttt{("generative AI" OR "large language model" OR LLM)}\\
\texttt{AND (comprehension OR "code understanding")}\\
\texttt{AND (code OR software)}
\end{quote}

The database-specific syntax, search dates, filters, and all five ScienceDirect strings are archived in the supplementary search log. The Scopus result export was further restricted to articles and conference papers in English; records outside the protocol's 2021--2026 year window were removed before merging. IEEE Xplore was filtered to conferences and journals within 2021--2026; ScienceDirect was filtered to research articles, English, and 2021--2026. The planned ACM Digital Library search was not performed because access was unavailable. A preliminary overly broad comprehension query was refined for precision; both this change and the ACM omission were recorded with their possible recall implications. The search date limits actual coverage: the study does not claim to include publications appearing after 17 September 2026.

\subsection{Study Selection}
Seven predefined inclusion criteria (IC1--IC7) required empirical primary studies involving human programmers using GenAI coding assistance in realistic development or maintenance tasks, with directly measured code comprehension or a defensible close proxy, peer-reviewed publication, and usable English full text. Eight exclusion criteria (EC1--EC8) covered model-only evaluations, productivity-only or code-quality-only studies without comprehension outcomes, introductory exercises that did not represent realistic software engineering, secondary or non-peer-reviewed reports, tools unrelated to GenAI assistance, unavailable usable full text, and duplicate publications without additional evidence. An author-posted manuscript was used as full text only when linked to an independently identifiable peer-reviewed version.

Database exports were mapped into a common metadata schema with source provenance preserved. Duplicate candidates were flagged first by DOI, then by normalized title, author names, and year; potentially distinct versions of the same empirical study required additional inspection. Records were first assessed using titles and abstracts, followed by full-text eligibility checks of retained reports. A criteria-guided AI-assisted workflow recorded screening decisions, a principal exclusion reason, and supporting rationale. For full-text exclusions, one primary reason was retained per report. The protocol proposed a four-reviewer calibration pilot, independent double screening, and an inter-rater reliability assessment, but the available decision log does not demonstrate completed independent pilot ratings or a reliability statistic; those procedures are therefore not claimed as completed. Search and selection counts belong in Section~\ref{sec:results} and the PRISMA flow diagram rather than in this procedural description.

\subsection{Quality Assessment}
Study quality was appraised against eight review-specific questions: clarity of objectives (QA1), suitability of empirical methods (QA2), participant and programming-experience description (QA3), specification of the assistant and interaction (QA4), task realism (QA5), clarity of the comprehension measure (QA6), support for the conclusions from reported analysis (QA7), and discussion of limitations or threats to validity (QA8). Each item was scored Yes (1), Partially (0.5), or No (0). Not Applicable items were omitted from the maximum applicable score. The resulting percentage was defined as
\begin{equation}
Q=100\,\frac{\sum_{i=1}^{8}s_i}{N_{\mathrm{applicable}}}.
\label{eq:qa_score}
\end{equation}
The review protocol set a study-specific threshold of $Q\geq60\%$, finalized on 18 September 2026. This custom quality appraisal is not presented as a validated risk-of-bias instrument or as a universal threshold. AI-assisted preliminary scores were entered into the quality log; independent two-reviewer confirmation and resolution of differences remain unverified at the current writing stage. The appraisal outcomes will be reported in Section~\ref{sec:results} once confirmed.

\subsection{Data Extraction and Synthesis}
A structured extraction form linked each eligible report to its study design; participant group and experience; assistant tool, model, and interaction mode; development task and codebase context; comprehension instrument; comparator; reported findings and statistics; comprehension-related challenges; limitations; and source section or page. Each field was mapped to one or more research questions. The extraction distinguishes \emph{direct} comprehension evidence (e.g., code-behavior questions or explanation accuracy) from \emph{close proxies} (e.g., cognitive load or perceived understanding); task completion speed alone is not coded as comprehension improvement. AI-assisted draft extraction entries remain subject to verification against the primary reports.

The planned cross-study synthesis groups evidence by research question rather than narrating one paper at a time. For RQ1, effects will be coded with the reported comprehension measure and direction (improved, degraded, mixed, no statistically significant difference, or not directly compared), without equating non-significance with equivalence. RQ2 groups tools and support forms by reported interaction mode; RQ3 classifies measurement instruments, participant experience, study design, and codebase/task context; and RQ4 develops descriptive themes for challenges, proposed mechanisms, and consequences. Theme labels and any ambiguous or discordant coding decisions will be checked against cited evidence locations and documented before the final synthesis. A narrative/thematic synthesis is planned because outcome measures and study designs are heterogeneous; meta-analysis will be considered only if sufficiently comparable effect estimates exist.

The protocol additionally specifies one round of backward and forward snowballing, citation tracking through Scopus, author-based checking, and a reference-list completeness check. These supplementary searches and the final synthesis have not yet been recorded as completed. Their outcomes will be incorporated in the final selection flow and results when available.

\section{Results}\label{sec:results}
\todo{Report database retrieval counts, duplicate removal, screening and full-text eligibility in the final PRISMA flow, then present study characteristics, quality appraisal results, and RQ1--RQ4 synthesis.}

\section{Discussion}\label{sec:discussion}
\todo{To be written.}

\section{Threats to Validity}\label{sec:threats}
\todo{To be written.}

\section{Conclusion}\label{sec:conclusion}
\todo{To be written.}

\bibliographystyle{IEEEtran}
\bibliography{references}

\end{document}
